\setlength{\parindent}{0pt}
% ------------------------------------------------------------
\chapter{Preliminaries}
\label{chap:preliminaries}
% ------------------------------------------------------------
\section{Notation and Mathematical Background}
\label{sec:notation}

This chapter collects the background the rest of the thesis builds on and fixes
the notation for it. We use $\{0,1\}$ or $\mathbb{F}_2$ for the set of
bits. Tuples are set in boldface, $\mathbf{x} = (x_1, \dots, x_n)$, with $|\mathbf{x}|$
denoting their length. For a finite set $X$ we write $\mathcal{P}(X)$ for its power set and
$\binom{X}{k}$ for the family of its $k$-element subsets. Functions are total
unless we say otherwise, and we identify a function $\{0,1\}^n \to \{0,1\}$ with the
Boolean function it computes. Notation that is used only locally is introduced
where it appears.
  \oldsubsection{\texorpdfstring%
    {The Field $\mathbb{F}_2$ and Multilinear Boolean Polynomials}%
    {The Field F2 and Multilinear Boolean Polynomials}}
  \label{sec:field-f2}

  The field $\mathbb{F}_2$ has carrier $\{0,1\}$ with addition $\oplus$, the
  exclusive-or, and multiplication $\odot$, the logical-and. Addition is its own
  inverse, %Suha: Grammar check 
  so $x \oplus x = 0$, and multiplication is idempotent, $x \odot x = x$.

  Every function $f : \mathbb{F}_2^{\,n} \to \mathbb{F}_2$ is a polynomial over
  $\mathbb{F}_2$, and idempotence lets us take that polynomial to have no variable
  raised above the first power. So $f$ has a unique \emph{multilinear}
  representation,
  \[
    f(x_1, \dots, x_n) \;=\; \bigoplus_{S \subseteq [n]} c_S \bigodot_{i \in S} x_i,
    \qquad c_S \in \mathbb{F}_2,
  \]
  its \emph{algebraic normal form} (ANF). Uniqueness is what turns the ANF into a
  decision procedure for the questions we put to it. In particular, note that, $f$ depends on a
  variable $x_i$ (meaning that flipping $x_i$ can change the output) if and only if
  some monomial of the ANF contains $x_i$.

  \subsection{Distributions, Laws, and Independence}
  \label{sec:prob-independence}

  Probing security is a statement about distributions, so we fix the little background 
  we need here.

  A \emph{distribution} on a finite set $X$ is
  a function $\mu : X \to [0,1]$ with $\sum_{x \in X} \mu(x) = 1$. It extends to
  subsets by summing, $\mu(A) = \sum_{x \in A} \mu(x)$ for $A \subseteq X$, and in
  that form it is called a \emph{measure}, so a distribution is a measure of total weight
  $1$, and we use the two words interchangeably. The \emph{uniform} distribution
  assigns $\mu(x) = 1/|X|$ to every point, hence $\mu(A) = |A|/|X|$ to every
  subset. On a cube $X = \{0,1\}^{n}$ the uniform distribution
  is the \emph{product} of the uniform distribution on each of the $n$ coordinates,
  giving $\mu(x) = 2^{-n}$ for every $x \in X$.

  %Suha (suggestion): Since you are using \mathcal{L} for laws, you can start introducing this notation here.
  A random variable $Z$ into a finite set $X$ induces
  the distribution $x \mapsto \Pr[Z = x]$, its \emph{law}. Two random variables are
  \emph{identically distributed} when their laws agree on every value.

  The security notion we are after is the independence of a family of laws from a
  parameter, which we state below.

   %Suha (suggestion): Since s is a parameter of Z, why don't you use Z(s) instead of Z_s in the definition below? 
  \begin{definition}[Independence of a parameter]
  \label{def:independence}
  Let $Z$ be a random variable into $X$ whose distribution depends on a parameter
  $s$ ranging over a set $\Sigma$. Write $Z_s$ for $Z$ under the value $s$ and
  $\mu_s$ for the law of $Z_s$, that is $\mu_s(z) = \Pr[Z_s = z]$, so that $s$
  selects one distribution out of the family $(\mu_s)_{s \in \Sigma}$. We say $Z$ is
  \emph{independent of} $s$ when
  \[
    \mu_s \;=\; \mu_{s'} \qquad \text{for all } s, s' \in \Sigma.
  \]
  \end{definition}

  We prove that two laws are equal by carrying one onto the other
  with a map that does not disturb the distribution.

  \begin{definition}[Measure-preserving map]
  \label{def:measure-preserving}
  A map $\pi : X \to X$ on a finite set is \emph{measure-preserving} for a
  distribution $\mu$ when
  \[
    \mu\big(\pi^{-1}(A)\big) \;=\; \mu(A) \qquad \text{for every } A \subseteq X .
  \]
  \end{definition}

  If $Z$ has law $\mu$, then $\pi(Z)$, the random variable obtained by applying
  $\pi$ to each outcome, has law $A \mapsto \mu(\pi^{-1}(A))$, which is called the
  \emph{pushforward} of $\mu$ along $\pi$. The condition in the above definition says exactly that
  the pushforward is $\mu$ again.

  For the uniform distribution the condition becomes bijectivity. As $|\pi^{-1}(A)| = |A|$
  for every $A$ forces $\pi$ to be injective, and an injection of a finite set into
  itself is onto. Every measure-preserving map in this thesis is of this kind, so
  the phrase may be read as ``bijective'' throughout. The technique we use
  simply states that if $\pi$ preserves the uniform distribution on $X$ and $Z$ is
  uniform, then for any function $g$ the variables $g(Z)$ and $g(\pi(Z))$ are
  identically distributed.

  \subsection{Boolean Circuits and Gates}
  \label{sec:bool-circuits}

We now introduce the object the rest of the thesis is built on.
The definition is the standard digraph formulation of a circuit, specialised to
arithmetic over the two-element field. On top of that graph we then read off, for
each wire, the $\mathbb{F}_2$ polynomial it computes. So, the graph is the circuit,
the polynomial on the wire is the
\emph{syntax} our checker rewrites, and the function that polynomial denotes is
the \emph{semantics} every claim in the thesis is ultimately about.

\begin{definition}[Arithmetic circuit over $\mathbb{F}_2$]
\label{def:circuit}
Fix the field $\mathbb{F}_2 = \{0,1\}$ and a set $S = \{s_i\}$ of \emph{gate
functions}, each of some arity $a_i \in \mathbb{N}_{\geq 1}$, i.e.\ each $s_i$
maps $\mathbb{F}_2^{\,a_i}$ to $\mathbb{F}_2$. A \emph{circuit} is a triple
$C = (\mathcal{W}, \lambda, \mathrm{op})$ in which
\begin{itemize}
  \item $\mathcal{W}$ is a finite set whose elements we call the \emph{wires},
  the vertices of the graph, partitioned as
  %Suha: What is the \uplus operation? Is it different than the regular set union?
  \[
    \mathcal{W} \;=\;
      \underbrace{\mathcal{S} \uplus \mathcal{R} \uplus \mathcal{P}}_{\text{input wires } \mathcal{V}}
      \;\uplus\; \mathcal{G}
  \]
  into \emph{secret}, \emph{random} and \emph{public} inputs (jointly constituting the
  input wires $\mathcal{V}$, and each carrying a variable) together with the
  \emph{computation gates} $\mathcal{G}$,
  \item $\lambda$ labels each computation gate $w \in \mathcal{G}$ with a gate function
  $\lambda(w) \in S$,
  \item $\mathrm{op}$ assigns to each computation gate $w \in \mathclap{G}$ its \emph{operand
  tuple} $\mathrm{op}(w) = (u_1, \dots, u_a) \in \mathcal{W}^{a}$, where $a$ is
  the arity of $\lambda(w)$.
\end{itemize}
subject to the requirement that the digraph $(\mathcal{W}, E)$ with edge set
\[
  E \;=\; \big\{\, (u, w) \mid w \in \mathcal{G}, \ u \in \mathrm{op}(w) \,\big\}
\]
contain no directed cycle. The input wires $\mathcal{V}$ have
indegree $0$; a computation gate has indegree the arity of its label, and a vertex of
outdegree $0$ is an \emph{output gate}. For a computation gate we speak
interchangeably of the gate and the wire carrying its output.
\end{definition}

% Recording the operands as a tuple, rather than reading them off the incoming
% edges, is deliberate. It orders them, which a non-commutative gate function would
% need, and it lets a gate take the same wire twice, as in $w = u \odot u$; the edge
% set $E$ forgets both and is used only to state acyclicity.

Over $\mathbb{F}_2$ there is no difference between the Boolean and the arithmetic
view of such a circuit: exclusive-or is addition, logical-and is
multiplication, and the two together already express every function
$\mathbb{F}_2^{\,n} \to \mathbb{F}_2$, i.e. they are universal. So we fix
\[
  S = \{\,\oplus,\ \odot\,\},
\]
addition and multiplication in $\mathbb{F}_2$ (Section~\ref{sec:field-f2}), both
of arity $2$. Every computation gate then has indegree exactly $2$, with
$\mathrm{op}(w_i) = (w_j, w_k)$, and reads
\[
  w_i = w_j \oplus w_k \qquad \text{or} \qquad w_i = w_j \odot w_k .
\]
Throughout the thesis $\odot$ denotes multiplication in $\mathbb{F}_2$; we
might sometimes use $\cdot$ for the same operation.

Acyclicity is what makes a circuit computable. It lets us list the wires in a
topological order $w_1, \dots, w_m$ in which the two operands of every gate
appear before the gate itself, so that $j, k < i$.

  \subsubsection{The Syntax of a Wire}

  Reading the graph along that order turns each wire into a polynomial over
  $\mathbb{F}_2$ in the input variables. The polynomials in question are the
  \emph{expressions}
  \[
    e \;::=\; 0 \;\mid\; 1 \;\mid\; v \;\mid\; e_1 \oplus e_2 \;\mid\; e_1 \odot e_2,
    \qquad v \in \mathcal{V} = \mathcal{S} \uplus \mathcal{R} \uplus \mathcal{P},
  \]
  written as they are built in the graph.

  \begin{definition}[The expression of a wire]
  \label{def:wire-expr}
  For $w \in \mathcal{W}$ the expression $\wexpr{w}$ is defined by recursion along
  the topological order:
  \[
    \wexpr{w} =
    \begin{cases}
      v & \text{if } w \in \mathcal{V} \text{ denotes the variable } v, \\
      \wexpr{w_j} \star \wexpr{w_k}
        & \text{if } w \in \mathcal{G} \text{ and } \mathrm{op}(w) = (w_j, w_k),\ \star = \lambda(w).
    \end{cases}
  \]
  \end{definition}

  The recursion is well founded exactly because the graph is acyclic. So, the
  operands of $w_i$ are $w_j, w_k$ with $j, k < i$, hence the second clause always
  descends. What it produces mentions no wire and only the inputs. The sharing
  present in the graph is unfolded away.

  \subsubsection{The Semantics of a Wire}

  Take as an \emph{input assignment} a map $\eta : \mathcal{V} \to \mathbb{F}_2$. Since
  $\mathcal{V}$ is a disjoint union, we write it as a triple
  $\eta = (\mathbf{a}, \rho, \pi)$ with a secret assignment
  $\mathbf{a} \in \{0,1\}^{\mathcal{S}}$, a \emph{tape}
  $\rho \in \{0,1\}^{\mathcal{R}}$, and a public assignment
  $\pi \in \{0,1\}^{\mathcal{P}}$. The expression of a wire denotes a bit value under
  such an assignment, by structural recursion on the grammar:
  \[
    \sem{0}\eta = 0, \qquad
    \sem{1}\eta = 1, \qquad
    \sem{v}\eta = \eta(v),
  \]
  \[
    \sem{e_1 \oplus e_2}\eta = \sem{e_1}\eta \oplus \sem{e_2}\eta, \qquad
    \sem{e_1 \odot e_2}\eta = \sem{e_1}\eta \odot \sem{e_2}\eta ,
  \]
  where the two operators on the left of each of the last two clauses are syntax
  and the ones on the right are the field operations of
  Section~\ref{sec:field-f2}. This is the usual reading of a polynomial as the
  function it computes.

  \begin{definition}[The wire function]
  \label{def:wire-function}
  The \emph{wire function} of $w \in \mathcal{W}$ in the circuit $C$ is
  \[
    \sem{w}_C \;=\; \sem{\wexpr{w}} \;:\;
      \{0,1\}^{\mathcal{S}} \times \{0,1\}^{\mathcal{R}} \times \{0,1\}^{\mathcal{P}}
      \;\longrightarrow\; \mathbb{F}_2 .
  \]
  \end{definition}

  So a wire \emph{is}, semantically a Boolean function, and equivalently,
  by Section~\ref{sec:field-f2}, a multilinear $\mathbb{F}_2$ polynomial function,
  which eats a secret, a random and a public assignment and returns one bit.
  The wire functions are built on top of one another over the shared variables of $\mathcal{V}$,
  since a gate's function is its two operands' functions combined by $\oplus$ or
  $\odot$, with the input wires supplying the variables everything else is built
  from. Fixing $\eta$ and evaluating gate by gate along the topological order
  computes exactly $\sem{w}_C(\eta)$ for every $w$, so the polynomial reading and
  the operational one are essentially the same.

  For example, a gate $c = p \oplus r$ whose operand $p$ is itself the gate
  $s \odot m$ has $\wexpr{c} = (s \odot m) \oplus r$, a polynomial in the secret
  $s$ and the randoms $m, r$, denoting the function
  $(\mathbf{a}, \rho) \mapsto \big(\mathbf{a}(s) \odot \rho(m)\big) \oplus \rho(r)$.

  Since the publics are held fixed throughout the thesis we suppress $\pi$, and we
  follow the usual abuse of notation to denote with $w$ both for the wire and for
  the function it denotes (i.e., $w(\rho; \mathbf{a}) = \sem{w}_C(\mathbf{a}, \rho, \pi)$).
  When the secrets are fixed as well we simply write $w(\rho)$. And a tuple of wires is made up
  componentwise.

\begin{notation}
\label{not:fv-vars}
The \emph{free variables} $\mathrm{fv}(e) \subseteq \mathcal{V}$ of an expression
are those it mentions, defined by structural recursion:
$\mathrm{fv}(0) = \mathrm{fv}(1) = \emptyset$, $\mathrm{fv}(v) = \{v\}$, and
$\mathrm{fv}(e_1 \star e_2) = \mathrm{fv}(e_1) \cup \mathrm{fv}(e_2)$. For a wire
we write $\mathrm{fv}(w) = \mathrm{fv}(\wexpr{w})$, so that
$\mathrm{fv}(w) = \{w\}$ for an input wire and
$\mathrm{fv}(w_j) \cup \mathrm{fv}(w_k)$ for a gate with operands $w_j, w_k$. We
set
\[
  \mathrm{vars}(e) \;=\; \mathrm{fv}(e) \cap \mathcal{R} \;\subseteq\; \mathcal{R}
\]
for the random inputs $e$ mentions, extended to a tuple of wires trivially by
$\mathrm{vars}(\mathbf{w}) = \bigcup_i \mathrm{vars}(w_i)$, and we call $e$
\emph{secret-free} when $\mathrm{fv}(e) \cap \mathcal{S} = \emptyset$.
\end{notation}

% Both notions are syntactic, defined by reading off an expression without evaluating it, which
% is what makes them cheap. The price is that they over-approximate genuine
% dependence --- $r \oplus r$ mentions $r$ but denotes $0$ --- and the following
% one-directional lemma is all we ever need of them.

% \begin{lemma}[Mentioning over-approximates dependence]
% \label{lem:fv-dependence}
% Let $e$ be an expression and $v \in \mathcal{V}$ with $v \notin \mathrm{fv}(e)$.
% Then $\sem{e}\eta = \sem{e}\big(\eta[v \mapsto b]\big)$ for every input assignment
% $\eta$ and every $b \in \mathbb{F}_2$.
% \end{lemma}

% \begin{proof}
% Structural induction on $e$. A constant does not read $\eta$. An input variable
% $u$ satisfies $u \neq v$, since $v \notin \mathrm{fv}(e) = \{u\}$, so it evaluates
% to $\eta(u)$ either way. For $e_1 \star e_2$ the free variables of both operands
% are contained in $\mathrm{fv}(e)$ and so also omit $v$, and the hypothesis applies
% to each.
% \end{proof}

% A secret-free expression is therefore independent of the secrets in the sense of
% Section~\ref{sec:prob-independence}. That implication is what the whole method
% runs on: the checker rewrites a probe until it is syntactically secret-free and
% concludes that it is semantically secret-independent. The converse fails for the
% syntactic test, and where an exact answer is needed we appeal instead to the
% uniqueness of the algebraic normal form (Section~\ref{sec:field-f2}), under which
% absence of a variable from the multilinear form is equivalent to independence of
% it.

% Two conventions, fixed here once, are used throughout. First, a topological
% order realises the circuit as a \emph{straight-line program}: each wire is
% defined at most once, in terms only of wires already defined, so there is no
% branching and no recomputation of a gate under a different name. We pass freely
% between the digraph of Definition~\ref{def:circuit} and this sequential reading,
% whichever is clearer at the moment. Second, the restriction to $2$-ary gates is
% a modelling choice, not a loss of generality: an $n$-ary sum or product is
% always expressible as a chain of $2$-ary gates.

\section{Secret Sharing and Masking}
\label{sec:secret-sharing}

Physical implementations of cryptographic hardware can leak information. Power draw, electromagnetic
emanation, and timing all correlate with the values a device computes on, and an
attacker who measures them can recover a key that is out of cryptographic reach
\parencite{kocher1999differential}. \emph{Masking} is the standard algorithmic
countermeasure \parencite{chari1999towards}. Instead of computing on a secret
directly, one splits it into several \emph{shares} whose joint distribution hides
it, then rewrites the computation to run on the shares throughout and reconstruct
only at the very end. The promise is that an attacker who sees few enough
intermediate values learns nothing, because too few shares carry no information
about the secret.

  \subsection{Boolean Masking}
  \label{sec:boolean-masking}

  We use \emph{Boolean} masking, where shares combine by exclusive-or. An
  \emph{$n$-sharing} of a bit $x$ is a tuple
  $(x_0, \dots, x_{n-1})$ with $x_0 \oplus \cdots \oplus x_{n-1} = x$. To draw one,
  sample $x_1, \dots, x_{n-1}$ uniformly and independently and set
  $x_0 = x \oplus x_1 \oplus \cdots \oplus x_{n-1}$.

  Placing this in the circuit of Definition~\ref{def:circuit} fixes each variable's
  role. The shared bit is a secret input, $x \in \mathcal{S}$. The $n-1$ freshly
  sampled shares are random inputs, $x_1, \dots, x_{n-1} \in \mathcal{R}$, and are
  exactly the coordinates a tape $\rho$ assigns. The remaining share $x_0$ is not
  an input but a computation gate, the $\oplus$ of $x$ with the others, so
  $x_0 \in \mathcal{G}$ with $\mathrm{fv}(x_0) = \{x, x_1, \dots, x_{n-1}\}$ and
  $\mathrm{vars}(x_0) = \{x_1, \dots, x_{n-1}\}$.

  Leaving out any one share destroys all information about $x$. Drop $x_j$ for any
  $j$ and let $Z$ be the tuple of the remaining $n-1$ shares, indicating a random variable.
  Whatever the value of
  $x$, that tuple is uniform on $\mathbb{F}_2^{\,n-1}$, so the law of $Z$ is the
  same for $x = 0$ and for $x = 1$. That is, it is independent of $x$ in the sense of
  Definition~\ref{def:independence}. The full tuple is not, since its sum is
  $x$ itself.

  A masked circuit carries every value in shared form and replaces each operation
  by a \emph{gadget}, a sub-circuit computing a sharing of the result from sharings
  of its inputs. Linear operations are trivial, since $\oplus$ acts share by share.
  Multiplication is where the difficulty and the fresh randomness live, such a gadget
  must combine shares of both inputs without ever forming a value that
  depends on a secret. The masked \textsc{and} of Ishai, Sahai and Wagner
  \parencite{10.1007/978-3-540-45146-4_27} is the archetype, and the circuits we
  verify are built from gadgets of this kind.

\section{The ISW Probing Model}
\label{sec:isw-model}

For the security notion to be precise we need to state an adversary. Ishai, Sahai and Wagner
\parencite{10.1007/978-3-540-45146-4_27} gave the model the thesis works in: the
probing model. An adversary places a bounded number of probes on the wires
of a circuit and reads the values they carry, and the circuit is secure when no
choice of probes tells the adversary anything about the secrets. The model
abstracts a physical side-channel attacker into a clean combinatorial game, and it
is the setting in which masking admits proofs.

The connections between the physical side-channel attack and the probing model, however,
are not exposited in this thesis.


  \subsection{Threshold Probing and Observation Sets}
  \label{sec:t-threshold}

  Fix a circuit with secret inputs $\mathcal{S}$, random inputs $\mathcal{R}$, and
  public inputs $\mathcal{P}$ as in Section~\ref{sec:bool-circuits}. An
  adversary of order $d$ selects at most $d$ wires and learns
  their joint values. That selection is the object the whole thesis reasons about,
  so we name it and give its notation.

  \begin{definition}[Probe]
  \label{def:probe}
  A \emph{probe} of order $d$ is a tuple
  $\mathbf{w} = (w_1, \dots, w_k)$ of $k \le d$ distinct wires of the circuit. Under
  a secret assignment $\mathbf{a} \in \{0,1\}^{\mathcal{S}}$ and a tape
  $\rho \in \{0,1\}^{\mathcal{R}}$ it takes the joint value
  \[
    \mathbf{w}(\rho; \mathbf{a}) \;=\;
      \big(w_1(\rho; \mathbf{a}), \dots, w_k(\rho; \mathbf{a})\big)
      \;\in\; \{0,1\}^{k} ,
  \]
  each component being the wire function of Definition~\ref{def:wire-function}
  with the publics held fixed. We write $\mathcal{L}_{\mathbf{a}}(\mathbf{w})$ for
  the \emph{law} of the probe under $\mathbf{a}$, i.e., the distribution of
  $\mathbf{w}(\rho; \mathbf{a})$ over a uniform pick of $\rho$.
  \end{definition}

  We follow the convention of Section~\ref{sec:notation} in setting the tuple in
  boldface, and we use $\mathbf{w}$ for a probe throughout the thesis. The
  literature also calls a probe an \emph{observation set}, and calls a value
  $o \in \{0,1\}^{k}$ that the probe can take an \emph{observation}. We use the
  word observation as well, and say ``probe'' for the tuple.

  % Listing the wires as a tuple rather than a set is a convenience, not a
  % commitment: the components are distinct and reordering them merely permutes the
  % coordinates of $\mathbf{w}(\rho; \mathbf{a})$, so it does not change whether the
  % law depends on the secrets. We therefore pass freely between the two readings,
  % writing $w \in \mathbf{w}$ for membership and $|\mathbf{w}| = k$ for the size. The
  % public inputs are known and the secrets are held fixed, so the only randomness in
  % play is the internal randomness $\mathcal{R}$, sampled uniformly, and every wire
  % of a probe is read as a function of $\rho$ alone.

  \subsection{\NoCaseChange{$d$}-Probing Security}
  \label{sec:d-probing}

  \begin{definition}[$d$-probing security]
  \label{def:d-probing}
  A circuit is \emph{$d$-probing secure} if every probe $\mathbf{w}$ of order $d$
  is \emph{secret-independent}. So, its law, taken over uniform
  $\rho \in \{0,1\}^{\mathcal{R}}$, is independent of the secret assignment in the
  sense of Definition~\ref{def:independence},
  \[
    \mathcal{L}_{\mathbf{a}}(\mathbf{w}) = \mathcal{L}_{\mathbf{b}}(\mathbf{w})
    \qquad \text{for all } \mathbf{a}, \mathbf{b} \in \{0,1\}^{\mathcal{S}} .
  \]
  \end{definition}

  The order $d$ counts the wires the adversary is granted, and an $n$-share
  masking is usually meant to withstand $d = n-1$ of them. The definition
  quantifies over tuples of wires read from a single run of the circuit, rather than single indivdual wires,
  since a masking leaks only jointly. And it is a worst-case statement over all
  $\sum_{i \le d} \binom{|\mathcal{W}|}{i}$ probes on the circuit's wires
  $\mathcal{W}$. Verifying this definition for the circuits we defined is our main
  task.

  % The displayed condition is Definition~\ref{def:independence} with the secret
  % assignment as the parameter and $\mathcal{L}_{\mathbf{a}}(\mathbf{w})$ as the
  % family of laws, so the law of a probe is one and the same whatever the secrets
  % are and we cannot deduce anything about them from the values on the probe. The
  % verification literature calls this condition \emph{probabilistic
  % non-interference} \parencite{barthe2015verified}; we note the name for the reader
  % who meets it elsewhere, but we keep to the equality of laws above, which is the
  % only form of the condition the thesis uses. Every statement we prove is an
  % equality of two laws, and every proof establishes one.

  % To show that a probe is non-interfering, it is enough to
  % rewrite its wires into a form in which the secrets have cancelled, which is the
  % program logic Barthe et al.\ introduce and the engine of
  % Chapter~\ref{chap:design} mechanises.

  % I figured, because we do not treat composability in the thesis. We don't really
  % need to introduce it in the preliminaries chapter.
  % \subsection{Threshold and Strong Non-Interference}
  % \label{sec:sni}

  % Probing security, unfourtanetly, does not compose. A secure gadget can
  % leak, once its output feeds another, because probes in the second gadget might see
  % combinations the first one never accounted for. Barthe et al.
  % \parencite{barthe2016strong} isolate two stronger notions that do compose. A
  % gadget is \emph{non-interfering} (NI) when any $d$ probes on it can be simulated
  % from at most $d$ shares of each input, and \emph{strongly non-interfering} (SNI)
  % when the same holds under a sharper budget that counts only the probes on
  % internal wires, not those on the outputs. SNI is what lets one assemble a large
  % masked circuit from small verified gadgets while preserving the order. These
  % notions organise much of the field, though the thesis stays with plain probing
  % security and does not pursue composition, so we recall them mainly to place our
  % object among its neighbours.

\section{Randomness Substitution}
\label{sec:randomness-substitution}

This section isolates, as a single reusable fact, the probabilistic argument
our checker automates. Fix a circuit as in Section~\ref{sec:bool-circuits},
with random inputs $\mathcal{R}$, and recall that a \emph{randomness assignment},
or \emph{tape}, is a function $\rho : \mathcal{R} \to \{0,1\}$. The secrets and public
inputs are held fixed throughout, so every wire $w$ denotes a function $w(\rho)$ of
$\rho$ alone (Definition~\ref{def:wire-function}). A probe $\mathbf{w}$ is as in
Definition~\ref{def:probe}, a tuple of at most $d$ wires whose order does not
affect the joint distribution, and Definition~\ref{def:d-probing} asks that this
joint distribution, taken over the uniformly distributed $\rho$, not move as the
secrets vary.

The whole argument rests on one line of $\mathbb{F}_2$ arithmetic: for any fixed
bit $c$, the translation/shear $b \mapsto c \oplus b$ is a \emph{bijection} of $\{0,1\}$
that preserves the uniform distribution (i.e., it is its own inverse and merely swaps
the two outcomes). Everything below is this fact, applied to a single coordinate
of $\rho$ with every other coordinate fixed.

Substitution is the syntactic half of that fact, and we mention that first. For
$r \in \mathcal{R}$ and an expression $t$, write $r \mapsto t$ for the
\emph{substitution} sending $r$ to $t$, and $e[r \mapsto t]$ for its application
to an expression $e$, which replaces every occurrence of $r$ in $e$ by $t$. The
brackets mark the application and the arrow names the substitution being applied.
Application is by structural recursion on the grammar of
Section~\ref{sec:bool-circuits}:
\[
  r[r \mapsto t] = t, \qquad
  u[r \mapsto t] = u \quad (u \in \mathcal{V} \cup \{0,1\},\ u \neq r),
\]
\[
  (e_1 \star e_2)[r \mapsto t] = e_1[r \mapsto t] \star e_2[r \mapsto t]
  \qquad (\star \in \{\oplus, \odot\}) .
\]
The grammar has no binders, so the recursion needs no side conditions.

\begin{definition}[Randomness substitution]
\label{def:substitution}
Let $r \in \mathcal{R}$ and let $e$ be an expression with
$r \notin \mathrm{vars}(e)$. The \emph{substitution}
$\sigma = (r \mapsto e \oplus r)$ acts on any expression $e'$ by
$e'\sigma = e'[\,r \mapsto e \oplus r\,]$, hence on a wire $w$ through its expression,
\[
  w\sigma \;=\; \wexpr{w}\,[\,r \mapsto e \oplus r\,] ,
\]
and on a probe componentwise. A chain of substitutions is applied in turn, each to
the expression the previous one produced. The same $\sigma$ acts on randomness assignments,
through the map $\widehat{\sigma} : \{0,1\}^{\mathcal{R}} \to
\{0,1\}^{\mathcal{R}}$ whose image $\rho' = \widehat{\sigma}(\rho)$ shifts the
$r$-coordinate by $e(\rho)$ and leaves every other coordinate untouched:
\[
  \rho'(r) = e(\rho) \oplus \rho(r), \qquad
  \rho'(r') = \rho(r') \quad \text{for } r' \neq r .
\]
\end{definition}

Geometrically, $\widehat{\sigma}$ is a \emph{shear} of the cube
$\{0,1\}^{\mathcal{R}}$, simply, it fixes every coordinate but $r$ and offsets that one
by $e(\rho)$. On the $r$-coordinate alone it is the translation
$\rho(r) \mapsto e(\rho) \oplus \rho(r)$, which merely swaps the two outcomes when
$e(\rho) = 1$ and does nothing when $e(\rho) = 0$, so it leaves that coordinate
uniform. We call any such map, one coordinate translated by a context
depending on the others, and everything else left intact, a \emph{shear}, and use
the word throughout.

We pair the two actions under one name with the lemma below. The
syntactic rewrite $w \mapsto w\sigma$ and the reparametrisation
$\rho \mapsto \widehat{\sigma}(\rho)$ are one operation seen from two sides, and
$\widehat{\sigma}$ does not disturb the uniform distribution.

\begin{lemma}[Substitution]
\label{lem:substitution}
For a substitution $\sigma = (r \mapsto e \oplus r)$ as in
Definition~\ref{def:substitution}, the map $\widehat{\sigma}$ is an involution of
$\{0,1\}^{\mathcal{R}}$ preserving the uniform distribution
(Definition~\ref{def:measure-preserving}), and for every wire $w$,
\[
  (w\sigma)(\rho) \;=\; w\big(\widehat{\sigma}(\rho)\big)
  \qquad \text{for every randomness assignment } \rho .
\]
\end{lemma}

\begin{proof}
Fix every coordinate of $\rho$ other than $r$. Since $r \notin \mathrm{vars}(e)$,
no leaf of $e$ is $r$, so evaluating $e$ never reads the $r$-coordinate and
$e(\rho)$ is a function of these fixed
coordinates alone; it is therefore constant as $\rho(r)$ ranges over $\{0,1\}$. On this
fibre $\widehat{\sigma}$ acts by the translation
$\rho(r) \mapsto e(\rho) \oplus \rho(r)$, which is its own inverse, so
$\widehat{\sigma}$ is an involution and in particular a bijection of
$\{0,1\}^{\mathcal{R}}$ --- hence measure-preserving for the uniform distribution.
Concretely, it leaves each coordinate uniform, and the uniform distribution on the
cube is the product of those, so the cube stays uniform.

For the stated identity, replace $w$ by its expression $\wexpr{w}$
(Definition~\ref{def:wire-expr}), which mentions only input variables, and argue
by structural induction on it. For a constant both sides are trivially equal. For
$\wexpr{w} = r$ both sides equal $e(\rho) \oplus \rho(r)$, on the left by the
clause $r[r \mapsto e \oplus r] = e \oplus r$, and on the right by the definition of
$\widehat{\sigma}$. For an input variable $u \neq r$ both sides equal $u(\rho)$,
since substitution leaves $u$ alone and $\widehat{\sigma}$ fixes every coordinate
but $r$. For $\wexpr{w} = e_1 \oplus e_2$ or $e_1 \odot e_2$, apply the hypothesis
to $e_1$ and $e_2$ and use that substitution distributes over both operators while
$\oplus$ and $\odot$ are computed pointwise from their operands' values.
\end{proof}

As a small remark, realize that the above lemma generalizes exactly as expected to a
tuple of wires $\mathbf{w}$ by coordinatewise application.

\begin{corollary}[Single-occurrence discharge]
\label{cor:single-occurrence-discharge}
Suppose $r$ occurs exactly once in the expression of a probe $\mathbf{w}$, as an
operand of a sum $x = e \oplus r$ with
$r \notin \mathrm{vars}(e)$, and let $\sigma = (r \mapsto e \oplus r)$. Then
$\mathbf{w}\sigma$ is identically distributed to $\mathbf{w}$ over $\rho$
uniform, and every occurrence of $x$ in $\mathbf{w}\sigma$ has collapsed to the
bare input $r$.
\end{corollary}

\begin{proof}
Distributional equality is Lemma~\ref{lem:substitution} applied coordinatewise
to $\mathbf{w}$, using that $\widehat{\sigma}$ preserves the uniform
distribution on $\{0,1\}^{\mathcal{R}}$. For the second claim,
$x\sigma = e \oplus (e \oplus r) = r$ by $\mathbb{F}_2$-cancellation, and since $r$
occurs nowhere else by hypothesis, no other part of $\mathbf{w}$ is touched.
\end{proof}

\begin{remark}
Rewriting a probe by a substitution as above --- whether in the general form of
Lemma~\ref{lem:substitution} or the single-occurrence case of
Corollary~\ref{cor:single-occurrence-discharge} --- leaves its distribution
untouched. That is the key observation that derives the procedure
to decide the security of a probe, by one
substitution at a time, taking it to a form with no secret in sight. The
single-occurrence case is the one the \texttt{maskVerif} tool
\parencite{10.1007/978-3-030-29959-0_15} calls \emph{optimistic sampling}. A
derivation is nothing more than such substitutions chained together; how that
chain is searched for is taken up in the later chapters.
\end{remark}

\section{Counting Complexity Classes}
\label{sec:hardness-preview}


Before closing off this chapter, we mention two specific complexity classes that
will come up later in the thesis. Let $M$ be a nondeterministic polynomial-time
Turing machine and write $\#M(x)$ for the number of accepting computation paths of $M$
on input $x$; where $\mathrm{NP}$ asks only whether $\#M(x) > 0$, the classes we are
concerned with look at the number itself.

\begin{itemize}
  \item $\#\mathrm{P}$, introduced by Valiant \parencite{valiant1979complexity},
  is the class of \emph{functions} of the form $x \mapsto \#M(x)$. It is a class
  of counting problems rather than decision problems. Its canonical member is
  $\#\textsc{Sat}$, the number of satisfying assignments of a Boolean formula.

  \item $\mathrm{C}_{=}\mathrm{P}$, studied by Wagner
  \parencite{wagner1986complexity}, is the class of \emph{languages}
  $\{\, x \mid \#M(x) = g(x) \,\}$ for a polynomial-time computable $g$, specifying if the
  machine have \emph{exactly} a prescribed number of accepting paths. This is a
  decision class, and it is the natural class of a yes-or-no question whose answer
  rests on a comparison of two counts. Its canonical complete problem is
  \textsc{ExactSat}, deciding whether a formula has exactly a given number of
  satisfying assignments.
\end{itemize}

Neither class is known to sit inside the polynomial hierarchy $\mathrm{PH}$. Actually, both
sit at least above it. Writing $\mathrm{P}^{\#\mathrm{P}}$ for the decision
problems solvable in polynomial time with an oracle for a counting function, two
observations place these classes. A formula is unsatisfiable exactly when its count is $0$,
which is an instance of exact counting, so $\mathrm{coNP} \subseteq
\mathrm{C}_{=}\mathrm{P}$; and two oracle calls, a subtraction and a test against
zero settle any exact-counting question, so $\mathrm{C}_{=}\mathrm{P} \subseteq
\mathrm{P}^{\#\mathrm{P}}$. Now, Toda's theorem states that $\mathrm{PH}
\subseteq \mathrm{P}^{\#\mathrm{P}}$ \parencite{toda1991pp}, so a single counting
oracle already decides every level of the hierarchy. In short,
\[
  \mathrm{P} \;\subseteq\; \mathrm{NP},\ \mathrm{coNP}
  \;\subseteq\; \mathrm{PH} \;\subseteq\; \mathrm{P}^{\#\mathrm{P}},
  \qquad
  \mathrm{coNP} \;\subseteq\; \mathrm{C}_{=}\mathrm{P}
  \;\subseteq\; \mathrm{P}^{\#\mathrm{P}} ,
\]
hence a problem complete for either class is well beyond anything we should expect to
solve exactly and efficiently.

Chapter~\ref{chap:couplings} shows that the two problems underlying $d$-probing
security (computing the law of one probe, and deciding whether two secrets
leave it unchanged) fall into exactly these two classes.
